%!TEX root = ../../main.tex
\section{다른 패치와 지역: 점수 전용 외부 평가}
\label{sec:val-external}

외부 평가는 이중 단계이다. 먼저 동결 평가기의 $V_{\mathrm{pre}} \to W$ 품질을, 다음으로 동결 $q$와 동결 기준선의 $\SVI$ 예측을 같은 공통 유효 행에서 봉인된 특징 순서로 점수화한다. 어떤 모델도 다시 적합하지 않았고 외부 어댑터도 두지 않았다. 표 \ref{tab:ext-T}와 \ref{tab:ext-S}는 점값이며, 주 두 표본(KR과 NA1의 16.13)에 대한 구간은 보완 부트스트랩(표 \ref{tab:ext-bootstrap})에서 따로 계산하였다.

%% table 객체: tab:ext-T — 출처 RRX_EXTERNAL 요약표; Q_NEWV_FIT85_TRANSFER_16X
\begin{table}[htbp]
  \centering
  \caption{외부 점수 전용 평가, T 행 (공통 유효 행; 동결 $\Vhat$, $q$, $\PTflex$; 구간 없음).}
  \label{tab:ext-T}
  \small
  \begin{tabular}{@{}lrrrrrrr@{}}
    \toprule
    코호트 & $n$ & $V_{\mathrm{pre}}$ Brier & $\dBrier(q - \PTflex)$ & $\Delta\MCB$ & $\Delta\DSC$ & $q$ $\MCB$ & $q$ $\DSC$ \\
    \midrule
    KR 16.13 & \numTExtKR{} & \extTKRVpre{} & \textbf{\extTKRdBrier{}} & \extTKRdMCB{} & \extTKRdDSC{} & \extTKRqMCB{} & \extTKRqDSC{} \\
    NA1 16.13 & \numTExtNA{} & \extTNAVpre{} & \textbf{\extTNAdBrier{}} & \extTNAdMCB{} & \extTNAdDSC{} & \extTNAqMCB{} & \extTNAqDSC{} \\
    KR 16.15 & \numTExtKRb{} & \extTKRbVpre{} & \extTKRbdBrier{} & --- & --- & --- & --- \\
    KR 16.14 파일럿 & \numTExtPilot{} & \extTPilotVpre{} & \extTPilotdBrier{} & --- & --- & --- & --- \\
    \bottomrule
  \end{tabular}
\end{table}


\paragraph{T (KR/NA1 16.13).} 두 외부 표본에서 $\dBrier(q - \PTflex)$의 점추정치는 \extTKRdBrier{}와 \extTNAdBrier{}로 기준선에 불리한 방향이다. 판별 성분은 $q$가 $\PTflex$보다 크고, 오보정 성분은 그보다 더 커서 순 Brier 격차의 부호를 정한다. $V_{\mathrm{pre}}$ Brier $\approx 0.15$는 평가기의 경기 승패 품질이 외부 표본에서 유지됨을 보이며, 외부 $\dV$ 라벨의 타당성은 별도의 질문이다. 여기의 CORP는 점값의 진단이다. $\PTlin$ 기준선에 대한 동반 전이표(\extTKRdBrierLin{} / \extTNAdBrierLin{})는 다른 기준선이므로 다른 양이다. 소규모 코호트(16.15, 파일럿)는 완전성을 위해 보고하며 어느 결론에도 합산하지 않는다.

%% table 객체: tab:ext-S — 출처 SCALE_SPLIT_TvsS_RESULTS json table4_EXT_S; RRX_EXTERNAL_S
\begin{table}[!t]
  \centering
  \caption{외부 점수 전용 평가: 소규모 교전 코호트 S}
  \label{tab:ext-S}
  \small
  \begin{tabular}{@{}lrrrrrrr@{}}
    \toprule
    코호트 & $n$ & $V_{\mathrm{pre}}$ Brier & $\dBrier(\qS - \PTflexS)$ & $\Delta\MCB$ & $\Delta\DSC$ & $\qS$ $\MCB$ & $\qS$ $\DSC$ \\
    \midrule
    KR 16.13 & \numSExtKR{} & \extSKRVpre{} & \textbf{\extSKRdBrier{}} & \extSKRdMCB{} & \extSKRdDSC{} & \extSKRqMCB{} & \extSKRqDSC{} \\
    NA1 16.13 & \numSExtNA{} & \extSNAVpre{} & \textbf{\extSNAdBrier{}} & \extSNAdMCB{} & \extSNAdDSC{} & \extSNAqMCB{} & \extSNAqDSC{} \\
    KR 16.15 & \numSExtKRb{} & \extSKRbVpre{} & \extSKRbdBrier{} & \extSKRbdMCB{} & \extSKRbdDSC{} & \extSKRbqMCB{} & \extSKRbqDSC{} \\
    KR 16.14 파일럿 & \numSExtPilot{} & \extSPilotVpre{} & \extSPilotdBrier{} & \extSPilotdMCB{} & \extSPilotdDSC{} & \extSPilotqMCB{} & \extSPilotqDSC{} \\
    \bottomrule
  \end{tabular}
  \tabnote{공통 유효 행; 동결 $\qS$, $\PTflexS$의 점수; 항등 보정기; 점값(16.13 두 표본의 구간은 표 \ref{tab:ext-bootstrap}).}
\end{table}


\paragraph{S (KR/NA1 16.13).} 같은 표본의 S 행에서는 점추정치의 순서가 KR과 NA1 모두에서 $\qS$를 $\PTflexS$ 앞에 둔다(\extSKRdBrier{}, \extSNAdBrier{}). \pendingauthor{외부 T/S 병기 문장의 유지 여부 (A2; E3 구간이 추가된 상태에서 결정)}

\paragraph{구간 (보완 부트스트랩).} 동결 예측을 그대로 두고 경기를 군집으로 복원 추출하는 부트스트랩(\ebBoot{}회, 시드 \ebSeed{})으로 네 주 대비(KR/NA1 $\times$ T/S)의 경기 평균 $\dBrier$에 95\% 백분위 구간과 네 대비 가족의 Bonferroni 98.75\% 구간을 계산하였다(표 \ref{tab:ext-bootstrap}; 보완 실험, 15.16 노출 이후의 진단). 구간은 고정된 모델에 조건부이며 모델 적합의 불확실성을 포함하지 않는다. T에서는 두 표본의 점추정치가 기준선에 불리한 방향이었다. KR에서는 차이의 불확실성이 컸고(\ebTKRD{}, 95\% 구간 \ci{\ebTKRLo}{\ebTKRHi}이 0을 포함), NA1에서는 95\% 구간이 양수였다(\ebTNAD{}, \ci{\ebTNALo}{\ebTNAHi}). 다중비교를 고려한 98.75\% 구간은 T NA1(\ci{\ebTNALoB}{\ebTNAHiB})에서만 0을 제외한다. S에서는 두 표본의 95\% 구간이 음수이고(KR \ebSKRD{} \ci{\ebSKRLo}{\ebSKRHi}, NA1 \ebSNAD{} \ci{\ebSNALo}{\ebSNAHi}) 98.75\% 구간은 0을 포함한다(KR \ci{\ebSKRLoB}{\ebSKRHiB}, NA1 \ci{\ebSNALoB}{\ebSNAHiB}). 전체 요약으로는, 평가한 외부 표본에서 T의 우세는 관측되지 않았고 S의 우세는 95\% 수준에서 관측되었으며 가족 보정 수준에서는 불확실하다.

%% GENERATED by thesis/tools/gen_supp.py -- do not edit
\begin{table}[!t]
  \centering
  \caption{외부 표본의 짝 경기 군집 부트스트랩 (16.13)}
  \label{tab:ext-bootstrap}
  \footnotesize

  \begin{tabular}{@{}llrrrrrll@{}}
    \toprule
    코호트 & 표본 & $n$ & 경기 & Brier $q$ & Brier PT & $\Delta$Brier & 95\% 구간 & 98.75\% 구간 \\
    \midrule
    T & KR 16.13 & 5{,}202 & 3{,}859 & 0.2483 & 0.2457 & +0.00261 & [-0.00041, +0.00572] & [-0.00115, +0.00668] \\
    T & NA1 16.13 & 5{,}312 & 3{,}955 & 0.2534 & 0.2494 & +0.00398 & [+0.00086, +0.00705] & [+0.00009, +0.00784] \\
    S & KR 16.13 & 15{,}641 & 7{,}874 & 0.2494 & 0.2512 & -0.00180 & [-0.00328, -0.00029] & [-0.00373, +0.00011] \\
    S & NA1 16.13 & 16{,}100 & 7{,}952 & 0.2503 & 0.2523 & -0.00204 & [-0.00363, -0.00044] & [-0.00403, +0.00001] \\
    \bottomrule
  \end{tabular}
  \tabnote{동결 $\hat V$, $q$, $\mathrm{PT}_{\mathrm{flex}}$(S는 $q_S$, $\mathrm{PT}_{\mathrm{flex}}^S$)의 점수; 재학습·재보정 없음. Brier는 경기 가중, $\Delta$Brier는 경기 평균 차이($q - \mathrm{PT}$, 음수이면 $q$ 우세). 구간은 10{,}000회 경기 군집 백분위 부트스트랩(시드 7); 98.75\%는 네 대비 Bonferroni. 구간은 고정된 모델에 조건부이며 모델 적합 불확실성을 포함하지 않는다.}
\end{table}


T와 S 행을 모두 가진 공통 경기(KR \ebHKRn{}, NA1 \ebHNAn{})에서 경기별 손실차의 차이 $H = \mathrm{mean}_m(d_{mT} - d_{mS})$는 KR \ebHKR{} \ci{\ebHKRLo}{\ebHKRHi}, NA1 \ebHNA{} \ci{\ebHNALo}{\ebHNAHi}이다(표 \ref{tab:ext-common-h}). $H$는 공통 경기 모집단의 대비이며 전체 코호트의 $D_T - D_S$와 다른 대상이고, T와 S의 차이에 대한 기전 진술의 근거로는 쓰지 않는다(제 \ref{sec:disc-scope} 절). 외부 확률 어댑터(F1)는 이연되어 있다(부록 \ref{app:deferred}).

%% GENERATED by thesis/tools/gen_supp.py -- do not edit
\begin{table}[!t]
  \centering
  \caption{공통 경기에서의 T$-$S 손실차 대비 (16.13)}
  \label{tab:ext-common-h}
  \small

  \begin{tabular}{@{}lrrl@{}}
    \toprule
    표본 & 공통 경기 & $H$ & 95\% 구간 \\
    \midrule
    KR 16.13 & 3{,}181 & +0.00364 & [-0.00063, +0.00790] \\
    NA1 16.13 & 3{,}290 & +0.00525 & [+0.00094, +0.00960] \\
    \bottomrule
  \end{tabular}
  \tabnote{$H = \mathrm{mean}_m(d_{mT} - d_{mS})$, $d_{mc}$ = 경기 $m$·코호트 $c$의 행 평균 $[(q-y)^2 - (\mathrm{PT}-y)^2]$; T와 S 행을 모두 가진 경기만. 전체 코호트의 $D_T - D_S$와 다른 대상이며 기전 진술의 근거가 아니다.}
\end{table}

